\documentclass[12pt, a4paper]{article}
\usepackage[utf8]{inputenc}
\usepackage[T2A]{fontenc}
\usepackage[russian]{babel}
\usepackage{amsmath, amssymb}
\usepackage{geometry}
\usepackage{graphicx}
\usepackage{hyperref}
\usepackage{setspace}
\usepackage{indentfirst}
\usepackage{caption}

\geometry{left=3cm, right=1.5cm, top=2cm, bottom=2cm}
\onehalfspacing
\setlength{\parindent}{1.25cm}

\captionsetup{justification=centering, labelsep=endash}

\begin{document}

\begin{titlepage}
  \centering
  {\small НАИМЕНОВАНИЕ ОРГАНИЗАЦИИ}\\[2cm]
  \textbf{ОТЧЁТ}\\[0.5cm]
  {\large\textbf{О НАУЧНО-ИССЛЕДОВАТЕЛЬСКОЙ РАБОТЕ}}\\[0.5cm]
  {\large\textbf{Тема работы}}\\[0.5cm]
  (промежуточный / заключительный)\\[3cm]

  \begin{flushleft}
    \hspace{8cm} Руководитель НИР \\
    \hspace{8cm} \underline{\hspace{3cm}} Фамилия~И.\,О. \\
  \end{flushleft}
  \vfill
  {\small Санкт-Петербург \\ \the\year}
\end{titlepage}

\newpage
\section*{Реферат}
\addcontentsline{toc}{section}{Реферат}
Отчёт содержит: XX~с., XX~рис., XX~табл., XX~источников.

КЛЮЧЕВЫЕ СЛОВА: слово1, слово2, слово3.

Объектом исследования является\ldots

Цель работы ---\ldots

\newpage
\tableofcontents
\newpage

\section*{Определения, обозначения и сокращения}
\addcontentsline{toc}{section}{Определения, обозначения и сокращения}

\begin{description}
  \item[API] --- Application Programming Interface
  \item[ГОСТ] --- Государственный стандарт
\end{description}

\newpage
\section*{Введение}
\addcontentsline{toc}{section}{Введение}

Актуальность\ldots

\newpage
\section{Основная часть}

\subsection{Анализ}
Текст.

\subsection{Результаты}
Текст.

\newpage
\section*{Заключение}
\addcontentsline{toc}{section}{Заключение}

Выводы\ldots

\newpage
\begin{thebibliography}{9}
  \bibitem{gost} ГОСТ 7.32-2017. Отчёт о научно-исследовательской работе.
\end{thebibliography}

\end{document}
